\documentclass[portrait]{seminar}
\usepackage{colordvi,boldmath,semcolor,semlayer,rotate,epsf,fancybox,mathsym}
\UseLegacyTextSymbols

\slidewidth9.5in
\slideheight6.7in
%\def\slidetopmargin{\iflandscape.2in\else0.8in\fi}
\def\slidetopmargin{\iflandscape.3in\else0.8in\fi}
\def\slidebottommargin{\iflandscape1.0in\else0.4in\fi}
\def\slideleftmargin{\iflandscape.2in\else.5in\fi}
\def\sliderightmargin{\iflandscape1.0in\else.8in\fi}

\newcommand{\cH}{{\cal H}}

\slidesmag{4}
\centerslidesfalse
\slideframewidth0.5pt
%\slideframe{none}

%\portraitonly
%\landscapeonly
%\def\leftsliderotation#1{\rotate[l]{#1}}
%\sliderotation{left}
\rotateheaderstrue


\newcommand{\onemat}{{\bf 1}}

\pagestyle{empty}

\begin{document}

\begin{slide}
{$ $}
\vfill
\begin{center}
{\Large \blue Computational Power of Hamiltonians\\ 
in Quantum Computing}

\vfill Pawel Wocjan

{\small Institut f\"ur Algorithmen und Kognitive Systeme}

{\small Universit\"at Karlsruhe}
\end{center}
\vfill
\end{slide}

%%%%%%%%%%%%%%%%%%%%%%%%%%%%%%%%%%%%
%
% research topics 
%
%%%%%%%%%%%%%%%%%%%%%%%%%%%%%%%%%%%%

\begin{slide}
{\large {\bf Important research topics of quantum computer science}}
\vfill
\begin{enumerate}
\item development of abstract {\red computational models} that
determine the elementary transformations of a quantum computer 

(based on physical laws governing the considered quantum system),
\item classification of {\red computational problems} with
respect to their complexity within these models
\item development of {\red efficient quantum algorithms} to handle these
various types of computational problems
\end{enumerate}
\vfill
development of {\it computational models} is central 

\hspace{4cm} {\huge $\Downarrow$} 

\hspace{3cm} complexity of problems and 

\hspace{3cm} efficiency of algorithms
\vfill
\end{slide}


\begin{slide}
{\large {\bf Control theoretic model of quantum computers}}
\vfill
respects more the properties of real quantum systems than the
quantum circuit model
\vfill
{\blue elementary operations of the control theoretic model}:
\begin{itemize}
\item external control operations in control
group $K$
\[
\mbox{       unitary transformations}\quad
V
\]
\item Hamiltonian time evolution internal to the system
\[
\exp(-i H t)
\]
\vfill
\end{itemize}
\end{slide}

\begin{slide}
{\large {\bf Control theoretic model of quantum computers}}
\vfill
implement $U$ by sequences

perform $V_1$, wait $t_1$, perform $V_2$, wait $t_2$, $\ldots$, 
perform $V_N$, wait $t_N$

such that
\[
U=\exp(-i H t_N) V_N \cdots \exp(-i H t_2) V_2 \exp(-i H t_1) V_1
\]
\vfill
{\red complexity of $U$}
\begin{itemize}
\item number of control operations $N$ and
\item time $t_1+t_2+\ldots+t_N$
\end{itemize}
\vfill
\end{slide}




\begin{slide}
{\large {\bf Mutual simulation of time evolutions}} 
\vfill 
use control theoretic model to implement
\begin{itemize}
\item transformations for solving computational problems
\item {\red transformations representing the time-evolution of another 
quantum system} ($\rightarrow$ universal quantum simulator)
\end{itemize}
\vfill 
simulation cost defines ``computational power'' of
different systems 
\vfill 
determining the cost: difficult!  
\vfill 
but: solvable for {\blue infinitesimal} time intervals
\vfill
\end{slide}

\begin{slide}
{\large {\bf Mutual simulation of {\red Hamiltonians}}} 
\vfill 
mutual simulation of time evolutions for {\red infinitesimal} time
intervals
\medskip

simulated time evolution $\exp(-i\, \tilde{H}\,{\red \epsilon})$ is
reproduced with error $O({\red \epsilon^2})$ 
\vfill
\[
\tilde{H} =
t_1\, {\blue U_1^\dagger H U_1}\, +\, 
t_2\, {\blue U_2^\dagger H U_2}\, +\, 
\ldots +\,
t_N\, {\blue U_N^\dagger H U_N} 
\]
\hspace{5cm}
{\LARGE $\Downarrow$}
\begin{eqnarray*}
&& 
\exp(-i\, {\blue U_N^\dagger H U_N}\, t_N{\red \epsilon})\, \cdots\, 
\exp(-i\, {\blue U_2^\dagger H U_2}\, t_2{\red \epsilon})\,
\exp(-i\, {\blue U_1^\dagger H U_1}\, t_1{\red \epsilon})  \\
& & \\
& = & 
\exp(-i\, \tilde{H}\,{\red \epsilon}) +O({\red \epsilon^2})
\end{eqnarray*}
\vfill
\end{slide}



\begin{slide}
{\large {\bf Mutual simulation of Hamiltonians}} 

\vfill
{\blue $H$ can simulate $\tilde{H}$ with overhead $1$}, written
$\tilde{H}\le H$, and {\blue $N$ time steps} iff
\[
\tilde{H}=\sum_{j=1}^N p_j\, U_j^\dagger H U_j\,,
\]
$p_1,\ldots,p_N>0$ and $p_1+\ldots p_N=1$

$U_j\in K$ are control operations
\medskip

{\blue $\tilde{H}$ is a convex sum of $N$ unitary conjugates of $H$}
\vfill
$H$ can simulate $\tilde{H}$ with overhead {\red $\tau$} iff 
$\tilde{H} \le {\red \tau} H$ 
\vfill 
{\red time-optimal simulation} problem $\longrightarrow$
{\red convex optimization} problem 
\vfill
\end{slide}

\begin{slide}
{\large {\bf Composite systems and pair-interaction Hamiltonians}}
\vfill
system consisting of {\blue $n$ qudits}  
$\quad {\cal H}=\C^d\otimes\C^d\otimes\cdots\otimes\C^d$ 
\vfill
{\blue pair-interaction Hamiltonian}
\[
H = 
\sum_{k<l} \sum_{\alpha\beta} J_{kl;\alpha\beta} \,\,
\sigma_\alpha^{(k)}\sigma_\beta^{(l)}\,,
\]
$B:=\{\sigma_\alpha\mid\alpha=1,\ldots,d^2-1\}$ is a basis of
$su(d)$
\vfill 
{\blue control operations} are all {\red local} operations in
\[
K:=SU(d)\otimes SU(d)\otimes\cdots\otimes SU(d)
\]
\vfill
\end{slide}




\begin{slide}
{\bf Main results of the thesis}
\vfill
\begin{itemize}
\item {\red formal definition of mutual simulation of Hamiltonians}
\item {\red lower and upper bounds on complexity} of mutual simulation
of Hamiltonians
\item {\red optimal simulation algorithms} using methods of
representation theory of finite groups (irreducible projective
representations) design theory (orthogonal arrays, difference schemes,
Hadamard matrices, and spreads) and algebraic graph theory
\item {\red application of mutual simulation of Hamiltonians to
adiabatic quantum computing} for solving computationally hard problems
\item {\red new computational problems and their quantum complexity}
classes
\end{itemize}
\vfill
\end{slide}

\begin{slide}
{\bf Algebraic Graph Theory gives Lower Bounds on Complexity}
\vfill
represent {\blue pair-interaction Hamiltonian}
\[
H = 
\sum_{k<l} \sum_{\alpha\beta} J_{kl;\alpha\beta}\,\, 
\sigma_\alpha^{(k)}\sigma_\beta^{(l)}
\]
\quad\quad\quad\quad\quad\quad 
$B:=\{\sigma_\alpha\mid\alpha=1,\ldots,d^2-1\}$ is a ONB of $su(d)$
\vfill
as {\blue coupling matrix} (extension of adjacency matrix of a graph)
\[
J=\left(
\begin{array}{c|c|c|c|c}
0      & J_{12} & J_{13} & \cdots & J_{1n} \\ \hline
J_{21} & 0      & J_{23} & \cdots & J_{2n} \\ \hline
\vdots & \vdots &        &        &        \\ \hline
J_{n1} & J_{n2} & J_{n3} &        & 0 
\end{array}
\right)\in\R^{(d^2-1)\,n\times (d^2-1)\,n}
\]
\vfill
\end{slide}

\begin{slide}
{\bf Analogy: coupling matrix $\leftrightarrow$ adjacency matrix}
\vfill
consider dipole-dipole coupling between all nodes
\[
H=\sum_{k<l} w_{kl} (3\,\sigma_z^{(k)}\sigma_z^{(l)} - 
\sum_{\alpha=x,y,z} \sigma_\alpha^{(k)}\sigma_\alpha^{(l)})\,,
\quad w_{kl}\neq 0
\]
\vfill
coupling matrix can be expressed as
\[
J={\red W}\otimes {\blue C}\,,
\]
where $W=(w_{kl})$ and $C={\rm diag}(-1,-1,2)$
\vfill
{\red $W$ adjacency matrix} of a weighted graph

{\blue $C$ characterizes the type of coupling} between nodes
\vfill
\end{slide}


\begin{slide}
{\large {\bf Action of control operations on the coupling matrix $J$}}
\vfill
\begin{center}
{\blue Conjugation of the pair-interaction Hamiltonian $H$ by control
operations}

corresponds to

{\blue conjugation of $J$ by a direct sum of orthogonal matrices}
\end{center}
\[
\begin{array}{ccc}
H & \mapsto & (U_1\otimes U_2\otimes\cdots\otimes U_n)^\dagger\, H\, 
(U_1\otimes U_2\otimes\cdots\otimes U_n) \\
J & \mapsto & (O_1\oplus O_2\oplus\cdots\oplus O_n)\, J\,
(O_1\oplus O_2\oplus\cdots\oplus O_n)^T
\end{array}
\]
\vfill
\end{slide}

\begin{slide}
{\large {\bf Lower bounds on time overhead}}
\vfill 
{\blue majorization criterion}

$H$ can simulate $\tilde{H}$ with time overhead $\tau$
\[
\sum_{j=1}^N p_j\, O_j \,(\tau J)\, O_j^T =
\tilde{J}
\]
$\tilde{J}$ convex sum of conjugates of $\tau J$
$\Longrightarrow$
{\red Spec($\tilde{J}$) majorizes Spec($\tau J$)}
\vfill
\end{slide}

\begin{slide}
{\large {\bf Lower bounds on time steps}} 
\vfill
{\blue rank criterion} (special case)

$H$ can simulate $\tilde{H}$ with $N$ time steps
\[
\tau\sum_{j=1}^N 
p_j\, O_j 
\,(W\otimes C)\,
O_j^T =
\tilde{W}\otimes C
\]
adding the block diagonal matrix $R:=\tau\sum_{j=1}^N p_j O_j
\,(\onemat\otimes C)\,O_j^T$ on both sides gives

$\Longrightarrow$ {\red $N \ge \mbox{rank}(\tilde{W}\otimes C + R)/
\mbox{rank}((W+\onemat)\otimes C)$}

\vfill
\end{slide}

\begin{slide}
{\large {\bf Complexity of time reversal}}
\vfill
time reversal: $H \mapsto -H$ (refocusing, spin echoes)
\vfill
$J=W\otimes C$ coupling matrix of $n$ coupled qubits, $w_{kl}\neq 0$ for 
all $k,l$
\vfill
\begin{enumerate}
\item $C$ traceless

time step and time overhead are at most $2$
\item $C$ has positive and negative eigenvalues, but $tr(C)\neq 0$

time steps at least $n/3-1$

time overhead independent of $n$, depends on eigenvalues only

\item $C$ either positive or negative semidefinite

time steps at least $n-1$

time overhead at least $n-1$
\end{enumerate}
\vfill

\end{slide}



\begin{slide}
{\large {\bf Complexity of time reversal}}
\vfill
{\red dipole-dipole coupling} between all nodes (easy to invert)
\begin{eqnarray*}
H & = & \sum_{k<l} (3\,\sigma_z^{(k)}\sigma_z^{(l)} - 
\sum_{\alpha=x,y,z} \sigma_\alpha^{(k)}\sigma_\alpha^{(l)}) \\
J & = & K \otimes {\rm diag}(-1,-1,2)\,,
\end{eqnarray*}
where $K$ is the adjacency matrix of the complete graph 
\vfill 
lower bound on time overhead is $2$ 
\vfill 
this bound is tight 

(well known-scheme in NMR with time overhead $2$ and $2$ time steps) 
\vfill
\end{slide}


\begin{slide}
{\large {\bf Complexity of time reversal}}
\vfill
{\red strong scalar coupling} between all nodes (difficult to reverse)
\begin{eqnarray*}
H & = & \sum_{k<l}
(\sigma_x^{(k)}\sigma_x^{(l)} + \sigma_y^{(k)}\sigma_y^{(l)} + 
 \sigma_z^{(k)}\sigma_z^{(l)})\\
J & = & K \otimes \onemat_3
\end{eqnarray*}
\vfill
lower bound on the time overhead $n-1$

$\longrightarrow$ reversed time-evolution necessarily slower by the 
factor $(n-1)$

lower bound on time steps $n-1$
\vfill
upper bounds $cn-1$ 

(decoupling algorithms based on orthogonal arrays)
\vfill
\end{slide}

\begin{slide}
{\large {\bf Computational power of Hamiltonians}}
\vfill
\begin{eqnarray*}
H_1 & = & \sum_{k<l} w_{kl} 
(\sigma^{(k)}_x\sigma^{(l)}_x\, {\red +}\, 
 \sigma^{(k)}_y\sigma^{(l)}_y) \\
& & \\
H_2 & = & \sum_{k<l} w_{kl} 
(\sigma^{(k)}_x\sigma^{(l)}_x\, {\blue -}\, 
 \sigma^{(k)}_y\sigma^{(l)}_y)
\end{eqnarray*}
\vfill
$H_1$ cannot simulate $H_2$ with overhead less than $(n-1)$

$H_2$ can simulate $H_1$ with overhead $2$
\vfill
\end{slide}

\begin{slide}
{\bf Graph-theoretical bounds: Selective decoupling}
\vfill
\epsfxsize0.35\textwidth\epsfbox{selective_decoupling.eps}
\hspace{0.5cm}
\raisebox{25ex}{\begin{minipage}[t]{5.5cm} 
complete Ising-Hamiltonian:

\quad $\sigma_z\otimes\sigma_z$  between all qubits

\bigskip

control operations $\onemat$ and $\sigma_x$
\end{minipage}}

{\blue Combinatorial problem: switching off unwanted interactions}

are the other interactions necessarily weakened (time overhead $>1$)?

yes, for some graphs $\rightarrow$ lower bound on time overhead given
by eigenvalues of the adjacency matrix

upper bound given by {\red chromatic index} of interaction graph
\vfill
\end{slide}

\begin{slide}
{\bf Selective decoupling based on chromatic index}
\vfill
\unitlength0.02cm
\begin{picture}(50,80)
\put( 0, 0){\line( 1, 0){40}}
\put(40, 0){\line( 0, 1){40}}
\put( 0,40){\line( 0,-1){40}}
\put( 0, 0){\line( 1, 1){40}}
\put( 0,40){\line( 1, 2){20}}
\put(20,80){\line( 1,-2){20}}
\put( 0, 0){\circle*{5}}
\put(40, 0){\circle*{5}}
\put(40,40){\circle*{5}}
\put( 0,40){\circle*{5}}
\put(20,80){\circle*{5}}
\end{picture}
\begin{picture}(40,80)
\put(15,30){\makebox(0,0){$=$}}
\end{picture}
\begin{picture}(50,80)
%\put( 0, 0){\line( 1, 0){40}}
%\put(40, 0){\line( 0, 1){40}}
%\put( 0,40){\line( 0,-1){40}}
\put( 0, 0){\line( 1, 1){40}}
\put( 0,40){\line( 1, 2){20}}
%\put(20,80){\line( 1,-2){20}}
\put( 0, 0){\circle*{5}}
\put(40, 0){\circle*{5}}
\put(40,40){\circle*{5}}
\put( 0,40){\circle*{5}}
\put(20,80){\circle*{5}}
\end{picture}
\begin{picture}(40,80)
\put(15,30){\makebox(0,0){$+$}}
\end{picture}
\begin{picture}(50,80)
%\put( 0, 0){\line( 1, 0){40}}
\put(40, 0){\line( 0, 1){40}}
\put( 0,40){\line( 0,-1){40}}
%\put( 0, 0){\line( 1, 1){40}}
%\put( 0,40){\line( 1, 2){20}}
%\put(20,80){\line( 1,-2){20}}
\put( 0, 0){\circle*{5}}
\put(40, 0){\circle*{5}}
\put(40,40){\circle*{5}}
\put( 0,40){\circle*{5}}
\put(20,80){\circle*{5}}
\end{picture}
\begin{picture}(40,80)
\put(15,30){\makebox(0,0){$+$}}
\end{picture}
\begin{picture}(50,80)
\put( 0, 0){\line( 1, 0){40}}
%\put(40, 0){\line( 0, 1){40}}
%\put( 0,40){\line( 0,-1){40}}
%\put( 0, 0){\line( 1, 1){40}}
%\put( 0,40){\line( 1, 2){20}}
\put(20,80){\line( 1,-2){20}}
\put( 0, 0){\circle*{5}}
\put(40, 0){\circle*{5}}
\put(40,40){\circle*{5}}
\put( 0,40){\circle*{5}}
\put(20,80){\circle*{5}}
\end{picture}
\vfill
Hadamard matrices: orthogonal matrices with $\pm 1$ entries

\epsfxsize0.29\textwidth\epsfbox{chromatic_index.eps}\hspace{0.1cm}\raisebox{10em}{
\begin{minipage}[t]{8cm}
$+1$: conjugation by $\onemat$
\medskip

$-1$: conjugation by $\sigma_x$
($\sigma_x \sigma_z \sigma_x = -\sigma_z$)
\medskip

rows: different pairs of qubits or single qubit
\medskip

columns: time steps
\end{minipage}}
\end{slide}

\begin{slide}
{\bf Application: ``one-way quantum computer''}
\vfill
required: Ising-interactions on two-dimensional lattice

but: resource Hamiltonian couples all qubits

$\rightarrow$ selective decoupling
\vfill
\epsfxsize0.7\textwidth\epsfbox{rectangular_lattice.eps}

$l\times l$ lattice
\vfill
selective decoupling with time overhead $4$ (=chromatic index)

lower bound on time overhead
\[
-4\cos\left( \frac{\pi}{l+1}l \right)\approx 4
\]
\vfill
\end{slide}

\begin{slide}
{\bf Application: continuous complexity measure for quantum circuits}
\medskip

weighted depth:

continuous complexity measure

based on control theoretic model with complete Ising-Hamiltonian

\hspace{1cm}
\unitlength0.03cm
\begin{picture}(100,150)
% 1 wire
\put(0,10){\line(1,0){10}}
% 2 wire
\put(0,30){\line(1,0){10}}
% 3 wire
\put(0,50){\line(1,0){10}}
% 4 wire
\put(0,70){\line(1,0){10}}
% 5 wire
\put(0,90){\line(1,0){10}}
%
% depth 1
\put(10,2){\framebox(20,16){}}
%
\put(10,22){\framebox(20,36){}}
%
\put(10,62){\framebox(20,36){}}
%
%
% 1 wire
\put(30,10){\line(1,0){10}}
% 2 wire
\put(30,30){\line(1,0){10}}
% 3 wire
\put(30,50){\line(1,0){10}}
% 4 wire
\put(30,70){\line(1,0){10}}
% 5 wire
\put(30,90){\line(1,0){10}}
%
% depth 2
\put(40,2){\framebox(20,36){}}
%
\put(40,42){\framebox(20,36){}}
%
\put(40,90){\line(1,0){20}}
%
% 1 wire
\put(60,10){\line(1,0){10}}
% 2 wire
\put(60,30){\line(1,0){10}}
% 3 wire
\put(60,50){\line(1,0){10}}
% 4 wire
\put(60,70){\line(1,0){10}}
% 5 wire
\put(60,90){\line(1,0){10}}
%
% depth 3
\put(70,10){\line(1,0){20}}
%
\put(70,22){\framebox(20,36){}}
%
\put(70,62){\framebox(20,36){}}
%
% 1 wire
\put(90,10){\line(1,0){10}}
% 2 wire
\put(90,30){\line(1,0){10}}
% 3 wire
\put(90,50){\line(1,0){10}}
% 4 wire
\put(90,70){\line(1,0){10}}
% 5 wire
\put(90,90){\line(1,0){10}}
%
\put(10,100){\makebox(20,20){1}}
\put(40,100){\makebox(20,20){2}}
\put(70,100){\makebox(20,20){3}}
%
\put(30,-15){a.\ depth}
\end{picture}
\hspace{1cm}
\begin{picture}(130,150)
% 1 wire
\put(0,10){\line(1,0){10}}
% 2 wire
\put(0,30){\line(1,0){10}}
% 3 wire
\put(0,50){\line(1,0){10}}
% 4 wire
\put(0,70){\line(1,0){10}}
% 5 wire
\put(0,90){\line(1,0){10}}
%
% weighted depth 1
\put(10,2){\line(0,1){16}}
%
\put(10,22){\framebox(30,36){}}
%
\put(10,62){\framebox(10,36){}}
%
%
% 1 wire
\put(10,10){\line(1,0){40}}
% 2 wire
\put(40,30){\line(1,0){10}}
% 3 wire
\put(40,50){\line(1,0){10}}
% 4 wire
\put(20,70){\line(1,0){30}}
% 5 wire
\put(20,90){\line(1,0){30}}
%
% weighted depth 2
\put(50,2){\framebox(20,36){}}
%
\put(50,42){\framebox(10,36){}}
%
% 1 wire
\put(70,10){\line(1,0){10}}
% 1 wire
\put(70,30){\line(1,0){10}}
% 3 wire
\put(60,50){\line(1,0){20}}
% 4 wire
\put(60,70){\line(1,0){20}}
% 5 wire
\put(50,90){\line(1,0){30}}
%
% weigthed depth 3
%
\put(80,22){\framebox(40,36){}}
%
\put(80,62){\framebox(30,36){}}
%
% 1 wire
\put(80,10){\line(1,0){50}}
% 2 wire
\put(120,30){\line(1,0){10}}
% 3 wire
\put(120,50){\line(1,0){10}}
% 4 wire
\put(110,70){\line(1,0){20}}
% 5 wire
\put(110,90){\line(1,0){20}}
%
%
\put(5,110){\vector(1,0){120}}
\put(5,105){\line(0,1){10}}
\put(125,105){\line(0,1){10}}
\put(55,110){\makebox(20,20){$\alpha$}}
%
\put(20,-15){b.\ weighted depth}
\end{picture}

\end{slide}

\begin{slide}
{\bf Application of mutual simulation of Hamiltonians to adiabatic quantum
computing}
\vfill
To do: construct {\red ``physically reasonable''} Hamiltonians whose ground 
states encode solutions of the problem ``independent set''
\begin{itemize}
\item graph $G=(V,E)$, positive integer $v\le |V|$
\item {\large $\exists$} independent set with cardinality at least $v$
\vfill 
subset $V'\subseteq V$ with 
\begin{itemize}
\item $|V'|\ge v$ and 
\item no two vertices in $V'$ are joined by an edge in $E$
\end{itemize}
\end{itemize}
\vfill
``independent set'' remains NP-complete for {\blue cubic planar} graphs.
\vfill
\end{slide}

\begin{slide}
{\bf Theorem (Planar spin glass within a magnetic field)}
\vfill
$G=(V,E)$ cubic planar graph 

determining {\blue energy of ground states} of the Hamiltonian
\[
\sum_{(k,l)\in E} \sigma_z^{(k)}\sigma_z^{(l)} + \sum_{k\in V}
\sigma_z^{(k)}
\]
equivalent to determining {\blue maximal cardinality of independent
sets}

\epsfbox[20 25 160 130]{graph3.ps}
\raisebox{14ex}{{but \red not physically reasonable!}}
\end{slide}

\begin{slide}
{\bf Planar orthogonal embedding} of graph $G=(V,E)$ maps

\begin{itemize}
\item vertices $k\in V$ to lattice points $\Gamma(k)$  
\item edges $(k,l)\in E$ to paths such that 
\begin{itemize}
\item images of their endpoints $\Gamma(k)$ and $\Gamma(l)$ are 
connected and 
\item paths do not intersect
\end{itemize}
\end{itemize}
\hspace{2cm}
\epsfxsize0.35\textwidth\epsfbox{planar_orthogonal_embedding.eps}
\vfill
\end{slide}

\begin{slide}
{\bf Planar orthogonal Hamiltonians}
\vfill
\hspace{0.25cm}
\epsfxsize0.35\textwidth\epsfbox{planar_orthogonal_Hamiltonian.eps}
\hspace{0.5cm}
\raisebox{15ex}{\begin{minipage}[t]{5.5cm} 
\begin{tabular}{rl}
filled circles & no $\sigma_z$ \\
thick lines & ${\red -c}\sigma_z\otimes\sigma_z$ \\
circles &  $\sigma_z$ \\
thin lines & $\sigma_z\otimes\sigma_z$
\end{tabular}
\end{minipage}}

ground states correspond $1$-to-$1$ to grounds states of $H$ for $c=3$

gap between ground and first excited is the same for $c=9$
\vfill
can be simulated efficiently by Ising-interactions on lattice

$16$ time steps, $2(c+1)$ time overhead
\vfill
\end{slide}

\begin{slide}
{\bf Quantum complexity classes}
\medskip

determining minimal energy of 
\begin{itemize}
\item pair-interaction ($2$-local) Hamiltonians: NP-complete
\item $k$-local Hamiltonians: {\blue QMA}-complete for $k\ge 3$
\end{itemize}
\medskip

QMA/QCMA quantum extensions of NP
\vfill
\end{slide}

\begin{slide}
{\bf New complete problems for QMA/QCMA}
\medskip

{\red low-energy and low-complexity states of local Hamiltonians}
\vfill

{\red Non-identity check}

\medskip

{\red Non-identity check on basis states}
\medskip

$\Rightarrow$ impression of how difficult it is to decide if an 
implementation in the control-theoretic model is optimal
\medskip

$\rightarrow$ justification for considering infinitesimal time
intervals (mutual simulation of Hamiltonians) 
\vfill
\end{slide}


\begin{slide}
{\bf Mathematical tools}

{\blue Vector majorization}

let $x=(x_1,x_2,\ldots,x_d)\in\R^d$

define
$x^\downarrow=(x^\downarrow_1,x^\downarrow_2,\ldots,x^\downarrow_d)\in\R^d$
as the vector whose components are the same as of $x$, but ordered so that
\[
x^\downarrow_1\ge x^\downarrow_2\ge\ldots\ge x^\downarrow_d
\]

{\blue $x$ majorizes $y$}, written {\blue $x\preceq y$}, if
\[
\sum_{j=1}^k x^\downarrow_j \le \sum_{j=1}^k y^\downarrow_j
\]
for $k=1,\ldots,d-1$ and 
\[
\sum_{j=1}^d x^\downarrow_j = \sum_{j=1}^d y^\downarrow_j
\]
\end{slide}

\begin{slide}
{\bf Mathematical tools} 

{\blue Operator majorization}
\vfill
let $A$ and $B$ be two Hermitian matrices of size $d$

we say {\blue $A$ majorizes $B$}, written {\blue $A\preceq B$}, if
\[
\lambda(A)\preceq\lambda(B)\,,
\]
where $\lambda(A)$ and $\lambda(B)$ are the spectra (the vectors of
eigenvalues) of $A$ and $B$, respectively.
\vfill
{\bf Uhlmann's theorem} 

we have $A\preceq B$ iff
\[
A=\sum_{j=1}^N p_j U_j^\dagger B U_j\,,
\]
where $\{p_j\}$ is a probability distribution and $\{U_j\}$ is a set
of unitaries
\vfill
\end{slide}

\end{document}

